Q-learning estimates action values by bootstrapping from $\max_{a'} Q(s',a')$. When the estimates are noisy, the maximum of estimates is on average larger than the maximum of the true values, an instance of the optimizer's curse~\cite{smith2006optimizer}. The resulting overestimation can slow learning, mislead action selection, and in function-approximation settings compound across bootstrapping steps. Double Q-learning (van Hasselt, 2010) keeps two estimators and uses one to pick the maximising action and the other to evaluate it, which removes the positive bias but can introduce a negative one~\cite{zhang2017weighted,ren2021estimation}. The idea carries into deep RL as Double DQN~\cite{hasselt2015deep}, is part of Rainbow~\cite{hessel2017rainbow}, and, as clipped double Q-learning, of TD3~\cite{fujimoto2018addressing}.

The standard evidence for the benefit of decoupling in the tabular setting is a single small MDP from Sutton and Barto's textbook (Example 6.7), in which the first action taken from the start state leads either to a terminal state with zero reward or to a state with many actions whose rewards are zero-mean-ish noise with a slightly negative mean. Beyond this example, the questions a practitioner would ask are less well answered in one place: how large is the bias on MDPs without a hand-built trap, does reducing it translate into better action choices at a fixed training budget, and how much of the observed difference between Q-learning and double estimators is due to decoupling rather than to confounds such as each table of a double learner being updated half as often, or the estimates not having converged from their initial values?

We run a controlled tabular study with exact true values. Our contributions are modest and empirical: (i) a vectorised numpy implementation of four estimators with identical step size, exploration and budget, evaluated with 1,000 independent runs per seed; (ii) a measurement of the estimated-minus-true start-state value and of the suboptimal-action fraction on the trap MDP and on random 20-state MDPs; (iii) a set of ablations (warm start at the true values, step size, exploration rate, noise level, number of actions, and updating both tables on each sample) that separate the effect of decoupling from learning speed and from adaptive sampling. The headline result is mixed: the trap-MDP result reproduces, but on random MDPs the double estimator underestimates more than Q-learning and takes more suboptimal actions at our training budget. We state our hypotheses and tests in Section~\ref{sec:setup} and give a verdict for each in Section~\ref{sec:results}. This is a small CPU study; Section~\ref{sec:limitations} lists what it does not cover.
